\documentclass[../../main.tex]{subfiles}
\begin{document}

\begin{center}
  {\fontsize{48pt}{58pt}\selectfont \textbf{1990}\normalsize\textbf{年度}}
\end{center}

\begin{table}[htb]
  \centering
  \caption{出題テーマと難易度}
  \begin{tabular}{|c|p{9cm}|c|}
    \hline
    \textbf{問題番号} & \textbf{テーマ} & \textbf{難易度} \\
    \hline
    第1問 & 接する三次・二次曲線に囲まれた面積 & \\
    \hline
    第2問 & 素数の辺をもつ三角形と余弦定理 & \\
    \hline
    第3問 & 行列による一次変換と内積 & \\
    \hline
    第4問 & 同心円間に並ぶ円の半径と面積 & \\
    \hline
    第5問 & 正多角形上の確率変数の独立性 & \\
    \hline
    第6問 & 円に外接し楕円に内接する三角形 & \\
    \hline
  \end{tabular}
\end{table}

\begin{table}[htb]
  \centering
  \caption{解答}
  \renewcommand{\arraystretch}{1.4}
  \begin{tabular}{|c|c|p{9cm}|}
    \hline
    \textbf{問題番号} & \textbf{小問} & \textbf{答え} \\
    \hline
    第1問 & -- & $\dfrac{4}{3}$ \\
    \hline
    第2問 & -- & $a=b=c$（証明問題） \\
    \hline
    第3問 & (1) &  \\
    \cline{2-3}
    & (2) &  \\
    \hline
    第4問 & (1) & $0<r\le\dfrac{\sin(\pi/n)}{1+\sin(\pi/n)}$ \\
    \cline{2-3}
    & (2) & $r=\dfrac{2}{n+4}$ \\
    \hline
    第5問 & (1) & $N=5$では独立でなく，$N=6$では独立 \\
    \cline{2-3}
    & (2) & $N=3,6$ \\
    \hline
    第6問 & -- & $a=\dfrac{b}{b-1}$（$b>1$） \\
    \hline
  \end{tabular}
\end{table}

\paragraph{試験形態}
6問．試験時間は要確認．

\end{document}
